% =====================================================================================================
\section{Background and Notation}\label{sec:background}
Table~\ref{tab:notation} lists the symbols used in the rest of the paper. All arithmetic of this section is the bit-exact contract between the golden model and the RTL: the hardware computes exactly these integer formulas, and the tests compare every intermediate value.

\begin{table}[t]
\caption{Notation.}
\label{tab:notation}
\centering\small\setlength{\tabcolsep}{3pt}
\resizebox{\columnwidth}{!}{%
\begin{tabular}{ll}
\toprule
Symbol & Meaning \\
\midrule
$W, H$ & frame width and height (multiples of 16) \\
$M = WH/256$ & number of 16$\times$16 macroblocks (MCUs) \\
$v$, $s$, $\amp(v)$ & quantized coefficient, its size category, its appended field \\
$t$, $\sigma$, $e$ & frame counter (63 bits), session identifier (64 bits), key epoch \\
$N(\sigma, t, b)$ & 128-bit nonce, $b \in \{0,1\}$ selects keystream or tag \\
$K_{\mathrm{m}}, K_e, \Kfb$ & master key, video key of epoch $e$, feedback key \\
$\Kseg, K_\pi$ & per-frame segment key, E5 permutation key \\
$Z_s$ & keystream of restart segment $s$ \\
$T_t$, $P_t$ & tag of frame $t$, its 26-byte authenticated prefix \\
$\pi_t$ & E5 permutation of the $M$ segments of frame $t$ \\
$L$, $c_L$, $\Theta_L$ & number of lanes, cycles per pixel, throughput \\
$\xor$, $\cat$, $\ggg$ & XOR, concatenation, rotation to the right \\
\bottomrule
\end{tabular}}
\end{table}

\subsection{Baseline JPEG}
\textit{Colour and sampling.} Each pixel $(R, G, B)$ is converted with 8-bit fixed-point coefficients and clamped to $[0, 255]$:
\begin{equation}
\begin{aligned}
Y  &= \lfloor (77R + 150G + 29B + 128) / 2^8 \rfloor,\\
C_b &= \lfloor (-43R - 85G + 128B + 32896) / 2^8 \rfloor,\\
C_r &= \lfloor (128R - 107G - 21B + 32896) / 2^8 \rfloor,
\end{aligned}
\label{eq:ycc}
\end{equation}
and every 2$\times$2 group of chroma samples $a, b, c, d$ is averaged as $\lfloor (a + b + c + d + 2)/4 \rfloor$ (4:2:0). A 16$\times$16 macroblock (MCU) therefore yields six 8$\times$8 blocks, four of luma and one of each chroma component, i.e.\ 384 samples for 256 pixels, or 1.5 samples per pixel.

\textit{Transform.} The two-dimensional DCT of a block $X$ of level-shifted samples is
\begin{equation}
F(u, v) = \tfrac{1}{4} c(u) c(v) \sum_{y=0}^{7} \sum_{x=0}^{7} X(y, x)\, \phi_v(x)\, \phi_u(y)
\label{eq:dctexact}
\end{equation}
with $\phi_k(x) = \cos\frac{(2x+1)k\pi}{16}$, $c(0) = 1/\sqrt2$ and $c(k) = 1$ otherwise. The hardware uses the separable integer form with the 13-bit matrix $C(u, x) = \operatorname{round}\bigl(2^{12} c(u) \phi_u(x)\bigr)$ and two rounding passes,
\begin{equation}
\begin{aligned}
G(y, v) &= \bigl\lfloor \bigl(\textstyle\sum_{x} C(v, x)\,(X(y, x) {-} 128) + 2^{10}\bigr) / 2^{11} \bigr\rfloor,\\
F(u, v) &= \bigl\lfloor \bigl(\textstyle\sum_{y} C(u, y)\,G(y, v) + 2^{14}\bigr) / 2^{15} \bigr\rfloor,
\end{aligned}
\label{eq:dct}
\end{equation}
whose scale is $2^{13} \cdot 2^{13} / 2^{26} = 1$, so $F$ approximates the exact DCT; the first pass keeps two fractional bits, which keeps the result within one unit of the exact transform.

\textit{Quantization.} Division by the table entry $Q(u, v)$ is replaced by a 17-bit reciprocal $R(u, v) = \operatorname{round}(2^{16}/Q(u,v))$,
\begin{equation}
q = \operatorname{sgn}(F) \bigl\lfloor (|F| R + 2^{15}) / 2^{16} \bigr\rfloor,
\label{eq:quant}
\end{equation}
clamped to $[-1024, 1023]$ for DC and $[-1023, 1023]$ for AC. With $R = 2^{16}/Q + \delta$, $|\delta| \le \tfrac12$, the product deviates from $|F|/Q$ by $|F|\,|\delta|/2^{16} \le 2^{11}/2^{17} = 2^{-6}$, so~\eqref{eq:quant} equals the rounded quotient except when $|F|/Q$ lies within $2^{-6}$ of a half-integer. One DSP slice per lane replaces a divider.

\textit{Entropy coding.} The 64 coefficients are read in zig-zag order. The DC coefficient is coded as the difference $d = q_0 - q_0'$ to the previous block of the same component; each non-zero AC coefficient as a pair (run of preceding zeros, size category) followed by an appended field. For $v \ne 0$ the size category is
\begin{equation}
s(v) = \lfloor \log_2 |v| \rfloor + 1, \qquad 2^{s-1} \le |v| \le 2^s - 1,
\label{eq:size}
\end{equation}
($s \le 11$ for DC differences, $s \le 10$ for AC), and the appended field is the $s$-bit one's-complement code
\begin{equation}
\amp(v) = \begin{cases} v, & v > 0,\\ v + 2^s - 1, & v < 0. \end{cases}
\label{eq:amp}
\end{equation}
The pair goes into a Huffman code (we use the tables of Annex~K of T.81~\cite{t81,wallace1992}); a run of 16 zeros is the symbol ZRL and the rest of a block after its last non-zero coefficient is EOB. Two properties follow and our design rests on both. First, the appended field can be changed without changing the Huffman code or the stream length, as long as $s$ is kept. Second, the Huffman symbols, and with them the run lengths and size categories of every non-zero coefficient, stay visible to anyone: they carry the outline of the picture.

\textit{Restart markers.} A restart marker RST$_m$ ($\mathrm{0xFFD0} + m$, $m = 0 \ldots 7$ cyclically) ends a restart segment: the bit stream is padded to a byte with 1-bits, all DC predictors are reset to zero, and the next segment can be decoded without any knowledge of the previous one~\cite{t81}. A DRI segment in the header gives the number of MCUs per segment. Inside segments, every byte 0xFF is followed by a stuffed 0x00, so a marker cannot occur in coded data. Restart segments make both the parallel lanes (Section~\ref{sec:lanes}) and the block permutation of E5 (Section~\ref{sec:e5}) possible.

\subsection{Ascon-AEAD128}\label{sec:ascon}
Ascon~\cite{ascon}, standardized by NIST in SP~800-232~\cite{sp800232}, is built on a 320-bit permutation applied to a state $S = x_0 \cat x_1 \cat x_2 \cat x_3 \cat x_4$ of five 64-bit words. One round is $p = p_L \circ p_S \circ p_C$:
\begin{itemize}
\item constant addition $p_C$: $x_2 \leftarrow x_2 \xor c_i$ with $c_i = 16(15 - i) + i$ for round $i$ (12 rounds: $i = 0 \ldots 11$; 8 rounds: $i = 4 \ldots 11$);
\item substitution $p_S$: a 5-bit S-box applied to the 64 bit slices $(x_0[j], \ldots, x_4[j])$ in parallel. In bitsliced form, with $a_0 = x_0 \xor x_4$, $a_4 = x_4 \xor x_3$, $a_2 = x_2 \xor x_1$ and $(t_0, \ldots, t_4) = (a_0, x_1, a_2, x_3, a_4)$,
\begin{equation}
b_i = t_i \xor (\neg t_{i+1} \wedge t_{i+2}) \quad (\text{indices mod } 5),
\label{eq:chi}
\end{equation}
followed by $x_0' = b_0 \xor b_4$, $x_1' = b_1 \xor b_0$, $x_2' = \neg b_2$, $x_3' = b_3 \xor b_2$, $x_4' = b_4$;
\item linear layer $p_L$: $x_i \leftarrow \Sigma_i(x_i) = x_i \xor (x_i \ggg r_i) \xor (x_i \ggg r_i')$ with $(r_i, r_i') = (19, 28), (61, 39), (1, 6), (10, 17), (7, 41)$.
\end{itemize}
Ascon-AEAD128 takes a 128-bit key $K$ and nonce $N$. It initializes $S \leftarrow p^{12}(\mathit{IV} \cat K \cat N) \xor (0^{192} \cat K)$; absorbs the associated data in 128-bit blocks, each XORed into $x_0 \cat x_1$ and followed by $p^8$; flips a domain-separation bit in $x_4$; encrypts each 128-bit plaintext block $P_i$ as $C_i = P_i \xor (x_0 \cat x_1)$, sets $x_0 \cat x_1 \leftarrow C_i$ and applies $p^8$; and finalizes with $S \leftarrow p^{12}(S \xor (0^{128} \cat K \cat 0^{64}))$ and the tag $T = (x_3 \cat x_4) \xor K$. Padding appends a single 1 and zeros to the last partial block. Bytes are loaded little-endian. The standard claims 128-bit security for confidentiality and integrity in the nonce-respecting setting and requires a new key when its data or decryption-failure limits are reached~\cite{sp800232}.

We use the construction in two ways. As a \emph{keystream generator}: for an all-zero message the ciphertext blocks are the rate words $x_0 \cat x_1$ after each $p^8$, i.e.\ a stream of pseudorandom bits under $(K, N)$. As a \emph{MAC}: the tag of an empty message with the data to authenticate as associated data. Both need the same core; one round per clock cycle gives 128 keystream bits every nine cycles (one cycle to absorb, eight rounds).

\subsection{Leakage and image-encryption metrics}\label{sec:metrics}
Let $I$ be the original luma and $\hat I$ a picture an attacker reconstructs. \emph{Edge correlation}, our main measure of outline leakage, is the Pearson correlation between Sobel gradient magnitudes,
\begin{equation}
\rho_E = \frac{\operatorname{cov}\bigl(|\nabla I|, |\nabla \hat I|\bigr)}{\sigma_{|\nabla I|}\,\sigma_{|\nabla \hat I|}},
\label{eq:edge}
\end{equation}
computed over all pixels; $\rho_E = 0$ means no outline, and it is insensitive to brightness and contrast, which a learned attacker recovers from global statistics. We also report SSIM~\cite{ssim}. For comparability with image-encryption papers we compute, between two pictures $C_1$, $C_2$ of $N$ pixels, $\mathrm{NPCR} = \frac{100}{N} \sum_{i} [C_1(i) \ne C_2(i)]$ and $\mathrm{UACI} = \frac{100}{N} \sum_{i} |C_1(i) - C_2(i)| / 255$, the histogram entropy $-\sum_k p_k \log_2 p_k$ and the correlation of adjacent pixels. These metrics were designed for ciphers that turn the whole image into noise; a format-compliant scheme keeps the stream decodable by design, so we read them together with~\eqref{eq:edge}.
